\ssscp{\tsc{Sumba}}{03-03-03-01}
While the evidence for a \tsc{Havu-Dhao} node is compelling,
the evidence for a \tsc{Sumba} node is much weaker.
The only evidence from shared sound changes for 
\tsc{Sumba} we have been able to identify is 
shared *z > \it{r} in four out of six
words for which reflexes have been identified.
In two other words *z = \it{ʤ},
though reflexes are currently only known in Kambera,
Lewa, and Mangili (\tsc{East Sumba}),
and one word may show *z > {\0}.
\tsc{Havu-Dhao} have *z > \it{ʄ} in the three identified reflexes.
Examples of *z in \tsc{Sumba-Havu} are given in \trf{03-58}.

\begin{table}\tcp{\tsc{Sumba-Havu} *z}{03-58}
\begin{adjustbox}{width=\textwidth}
	\begin{threeparttable}\st{2pt}
	\begin{tabular}{llllllll}\lsptoprule
*gloss	&	path	&	rain	&	far	&	near	&	widely spaced	&	Job’s tears	&	bad\\
PMP	&	*\tbr{z}alan\su{†}	&	*qu\tbr{z}an	&	*\tbr{z}auq	&	*qa\tbr{z}ani	&	*\tbr{z}araŋ	&	*\tbr{z}əlay	&	*\tbr{z}aqat\su{Δ}\\	\midrule
Kodi	&	{\hr}la\tbr{r}a	&	{\hq}u\tbr{r}ːa	&	{\hr}mba\tbr{r}oyo	&	\cg{(tukːa)}	&		&		&	\hp{*z}akic-o\\
Weyewa	&	{\hr}la\tbr{r}a	&	{\hq}u\tbr{r}ːa	&	{\hr}ma\tbr{r}ːo	&		&		&		&	\hp{*z}akut-a\\
Loli	&	{\hr}la\tbr{r}a	&	{\hq}u\tbr{r}ːa	&	{\hr}ma\tbr{r}ːo	&	\cg{(tukːe)}	&		&		&	\\
Laboya	&	{\hr}laː\tbr{r}a	&	{\hq}u\tbr{r}a	&	{\hr}ma\tbr{r}au	&	\cg{(daheka)}	&		&		&	\\
Wanukaka	&	{\hr}la\tbr{r}a	&	{\hq}u\tbr{r}a, u\tbr{r}aŋ&	{\hr}ma\tbr{r}ou	&	{\hr}ma\tbr{r}ani	&		&		&	\\
Mamboru	&	{\hr}la\tbr{r}a	&	{\hq}u\tbr{r}a	&	\cg{(malera)}	&	{\hr}ma\tbr{r}ani	&		&		&	\hp{*z}akat-a\\
Anakalang	&	{\hr}la\tbr{r}a	&	{\hq}u\tbr{r}aŋ-u	&	{\hr}ma\tbr{r}au	&		&		&		&	\hp{*z}akat-u\\
Lewa	&	{\hr}la\tbr{r}a	&	{\hq}u\tbr{r}aŋ-u	&	{\hr}\tbr{r}au	&	{\hr}maranu	&	{\hr}\tbr{ʤ}araŋ-u	&	{\hr}\tbr{ʤ}olu	&	\hp{*z}akat-u\\
Kambera	&	{\hr}la\tbr{r}a	&	{\hq}u\tbr{r}aŋ-u	&	{\hr}\tbr{r}au, ma\tbr{r}au&	{\hr}ka\tbr{r}eni	&	{\hr}\tbr{ʤ}araŋ-u	&	{\hr}\tbr{ʤ}uli\su{‡}&	\hp{*z}akat-u\\
Mangili	&	{\hr}la\tbr{r}a	&	{\hq}u\tbr{r}aŋ-u	&	{\hr}\tbr{r}au	&	{\hr}ka\tbr{r}eni	&	{\hr}\tbr{ʤ}araŋ-u	&	{\hr}\tbr{ʤ}uli	&	\hp{*z}akat-u\\	\midrule
Havu	&	{\hr}ru\tbr{ʄ}ara	&	{\hq}ə\tbr{ʄ}i	&	{\hr}\tbr{ʄ}əu	&	\cg{(umu)}	&		&		&	\hp{z}\cg{(apa)}\\
Dhao	&	{\hr}\tbr{ʄ}ara	&	{\hq}ə\tbr{ʄ}i	&	{\hr}ka\tbr{ʄ}əu	&	\cg{(dətu)}	&		&		&	\hp{z}\cg{(b͡βelu)}\\
		\lspbottomrule
	\end{tabular}
		\begin{tablenotes}\tnsize
			\item[†]	{Reflexes of *zalan have metathesis of the initial
								and medial consonants in Sumba.}
			\item[‡]	{\citet[71]{Onvlee1984} gives Kambera \it{ʤuli} as ‘small kind
								of grain, the pit of which is used as a bead’.}
			\item[Δ]	{If the words given here are reflexes of *zaqat,
								they also show irregular *q > \it{k}
								and irregular changes of the final vowel in Kodi and Weyewa.}
		\end{tablenotes}
	\end{threeparttable}
\end{adjustbox}
\end{table}

The change of *z > \it{r} in the four words given in \trf{03-58}
results in a partial merger of *R,
*j, and *z in \tsc{Sumba} and groups
these languages together apart from \tsc{Havu-Dhao}.
\citet{Gasser2014b} also finds lexical
support for a distinct \tsc{Sumba} group.

Within \tsc{Sumba} we can identify a number of groupings on the basis of lexical similarity.
The rate of lexical similarity among languages of Sumba is presented
in \trf{03-59}, based on \citet{Asplund2010} who calculated this with a modified
Swadesh word list of 210 items.
Languages are reordered (roughly from east to west)
and labels changed to match the language/dialect names in \frf{03-01}.
On the basis of lexical similarity five rough groupings emerge:
\tsc{East Sumba}, \tsc{Central Sumba}, \tsc{West Sumba}, Laboya, and \tsc{Kodi-Gaura}.

\begin{table}
\caption[\tsc{Sumba} lexical similarity (adapted from Asplund 2010)]
{\tsc{Sumba} lexical similarity (adapted from \citealt{Asplund2010})}
\label{tab:03-59}
\begin{adjustbox}{width=\textwidth}
	\st{3pt}\begin{tabular}{cccccccccccccccccccccccc}\hhline{--------------~~~~~~~~~~}
\mc{3}{|l}{Mangili}		&	\mc{11}{r|}{\tsc{East Sumba}}	&&&&&&&&&\\	\hhline{~------------~~~~~~~~~~~}
\mc{1}{|c|}{87 {\SR}}	&	\mc{6}{l}{Tabundung-Karera}		&	\mc{6}{r|}{\tsc{Nuclear E. Sumba}}	&	\mc{1}{c|}{}	&&&&&&&&&	\\	
\mc{1}{|c|}{84 {\SR}}	&	90 {\SR}&	\mc{11}{l|}{Kambera}		&\mc{1}{c|}{}&&&&&&&&&&		\\	
\mc{1}{|c|}{78 {\SY}}	&	83 {\SR}&	85 {\SR}&	\mc{10}{l|}{Kanatang}			&\mc{1}{c|}{}&&&&&&&&&&		\\	
\mc{1}{|c|}{77 {\SY}}	&	82 {\SR}&	84 {\SR}&	89 {\SR}&	\mc{9}{l|}{Prai Paha}			&\mc{1}{c|}{}&&&&&&&&&&	\\	
\mc{1}{|c|}{76 {\SY}}	&	83 {\SR}&	83 {\SR}&	85 {\SR}&	89 {\SR}&	\mc{8}{l|}{Lewa}			&\mc{1}{c|}{}&&&&&&&&&&		\\	
\mc{1}{|c|}{77 {\SY}}	&	80 {\SR}&	82 {\SR}&	90 {\SR}&	89 {\SR}&	84 {\SR}&	\mc{7}{l|}{Hahar}			&\mc{1}{c|}{}&&&&&&&&&&		\\	
\mc{1}{|c|}{75 {\SY}}	&	78 {\SY}&	81 {\SR}&	87 {\SR}&	84 {\SR}&	83 {\SR}&	90 {\SR}&	\mc{6}{l|}{Umbu Ratu Nggai}		&\mc{1}{c|}{}&&&&&&&&&		\\	\hhline{~------------~~~~~~~~~~~}
\mc{1}{|c}{69 \cgr}		&	71 {\SY}&	74 {\SY}&	79 {\SY}&	76 {\SY}&	75 {\SY}&	79 {\SY}&	85 {\SR}&	\mc{6}{l|}{Praiwatana}	&&&&&&&&&	\\	
\mc{1}{|c}{65 \cgr}		&	69 \cgr	&	69 \cgr	&	72 {\SY}&	70 {\SY}&	72 {\SY}&	73 {\SY}&	79 {\SY}&	84 {\SR}&	\mc{5}{l|}{Maderi}	&&&&&&&&&	\\	\hhline{------------------~~~~~~}
62 \cgr								&	63 \cgr	&	66 \cgr	&	70 {\SY}&	67 \cgr	&	67 \cgr	&	73 {\SY}&	74 {\SY}&	77 {\SY}&	74 {\SY}&	\mc{3}{|l}{Mamboru}	&\mc{5}{r|}{\tsc{Central Sumba}}	\\	
62 \cgr								&	66 \cgr	&	63 \cgr	&	67 \cgr	&	66 \cgr	&	66 \cgr	&	68 \cgr	&	69 \cgr	&	69 \cgr	&	73 {\SY}&	\mc{1}{|c}{79 {\SY}}&	\mc{7}{l|}{Anakalang}								\\	
55										&	56			&	56			&	59			&	60 \cgr	&	60 \cgr	&	60 \cgr	&	61 \cgr	&	69 \cgr	&	63 \cgr	&	\mc{1}{|c}{70 {\SY}}&	76 {\SY}	&	\mc{6}{l|}{Pondok}											\\	
58										&	63 \cgr	&	61 \cgr	&	62 \cgr	&	62 \cgr	&	62 \cgr	&	60 \cgr	&	62 \cgr	&	58			&	63 \cgr	&	\mc{1}{|c}{66 \cgr}	&	76 {\SY}	&	65 \cgr	&	\mc{5}{l|}{Wanukaka}							\\	
47										&	50			&	49			&	52			&	53			&	53			&	54			&	56			&	61 \cgr	&	60 \cgr	&	\mc{1}{|c}{62 \cgr}	&	72 {\SY}	&	76 {\SY}&	60 \cgr	&	\mc{4}{l|}{Baliledo}				\\	\hhline{~~~~~~~~~~--------------}
51										&	53			&	53			&	54			&	53			&	51			&	55			&	57			&	59			&	56			&	56									&	60 \cgr		&	61 \cgr	&	64 \cgr	&	61 \cgr	&	\mc{2}{|l}{Loli}			&		\mc{7}{r|}{\tsc{West Sumba}}	\\	
47										&	49			&	49			&	52			&	51			&	48			&	52			&	53			&	56			&	53			&	53									&	54				&	56			&	61 \cgr	&	58			&	\mc{1}{|c}{80 {\SR}}&	\mc{8}{l|}{Tana Righu}	\\	
53										&	48			&	47			&	49			&	49			&	47			&	49			&	51			&	54			&	48			&	51									&	52				&	53			&	64 \cgr	&	58			&	\mc{1}{|c}{74 {\SY}}&	80 {\SR}	&	\mc{7}{l|}{Waibangga}	\\	
43										&	46			&	47			&	50			&	49			&	45			&	51			&	50			&	53			&	49			&	51									&	51				&	53			&	56			&	54			&	\mc{1}{|c}{78 {\SY}}&	73 {\SY}	&	81 {\SR}&	\mc{6}{l|}{South Weyewa}	\\	
44										&	46			&	44			&	48			&	47			&	44			&	49			&	49			&	52			&	49			&	49									&	51				&	54			&	57			&	53			&	\mc{1}{|c}{79 {\SY}}&	84 {\SR}	&	72 {\SY}&	84 {\SR}&	\mc{5}{l|}{Weyewa}	\\	
43										&	47			&	47			&	49			&	48			&	46			&	49			&	51			&	53			&	49			&	51									&	51				&	53			&	51			&	54			&	\mc{1}{|c}{74 {\SY}}&	76 {\SY}	&	69 \cgr	&	82 {\SR}&	75 {\SY}&	\mc{4}{l|}{Bukambero}	\\	\hhline{~~~~~~~~~~~~~~~---------}
42										&	44			&	46			&	46			&	46			&	46			&	47			&	49			&	52			&	47			&	50									&	53				&	58			&	51			&	57			&	65 \cgr							&	64 \cgr		&	69 \cgr	&	63 \cgr	&	61 \cgr	&	61 \cgr	&	\mc{3}{l}{Laboya}	\\	\hhline{~~~~~~~~~~~~~~~~~~~~~~--}
48										&	49			&	50			&	53			&	51			&	49			&	53			&	54			&	54			&	51			&	53									&	56				&	60 \cgr	&	51			&	52			&	62 \cgr							&	62 \cgr		&	60 \cgr	&	62 \cgr	&	59			&	62 \cgr	&	63 \cgr	&	\mc{2}{|l|}{Gaura}	\\	
46										&	47			&	48			&	49			&	48			&	47			&	50			&	52			&	52			&	49			&	50									&	51				&	54			&	46			&	48			&	57									&	57				&	53			&	60 \cgr	&	54			&	67 \cgr	&	54			&	\mc{1}{|c}{74 {\SY}}&	\mc{1}{l|}{Kodi}	\\	\hhline{~~~~~~~~~~~~~~~~~~~~~~--}
\lspbottomrule
	\end{tabular}
\end{adjustbox}
\end{table}

Exclusively shared phonological innovations supporting these
lexical groupings have not been identified.
All members of \tsc{Central Sumba},
except Mamboru, have denasalisation of prenasalised stops: **mb
> \it{b}, **nd > \it{d}, and **ŋg > \it{g}
and this could be used to further group these four languages.
However, this sound change also occurs in Loli and Gaura.

Similarly, as discussed further in \srf{03-03-03-02},
many languages of western Sumba
have lengthening of certain word medial consonants.
%This is known to occur to some extent in Wanukaka,
%Weyewa, and Kodi and there is historical evidence that it occurred
%in Laboya (see \srf{03-03-01}).
However, the full extent of consonant lengthening is not known,
partly due to transcriptional and/or analytical differences.
This may provide evidence for a node containing \tsc{West-Sumba}, Laboya, and \tsc{Gaura-Kodi}.
However, we do not currently use it for subgrouping because of
the lack of sufficient data on its distribution.
%as well as the possibility that it was innovated in \tsc{Sumba-Havu}
%but lost in languages of eastern Sumba.

While several languages of Sumba,
as well as Havu, have *s > \it{h},
this change does not occur in Dhao nor in all languages of Sumba.
There is evidence that this shift happened in historical
times in both Kambera \citep[12]{Klamer1998a} and Havu \citep[3]{WalkerA1982}.
